% Reusable primitives for transformer-tensor-viz.
% Requires TikZ; callers may override colors before inputting this file.
\usetikzlibrary{arrows.meta,backgrounds,calc,matrix,positioning}

\providecolor{ttArrowBlue}{HTML}{82B5D5}
\providecolor{ttOutlineBlue}{HTML}{7892A3}
\providecolor{ttMutedText}{HTML}{65737C}
\providecolor{ttMaskCell}{HTML}{E3E6E8}

\tikzset{
  tt flow/.style={draw=ttArrowBlue,line width=0.8pt,-{Stealth[length=5pt,width=5pt]}},
  tt share/.style={draw=ttArrowBlue!88,dashed,line width=0.72pt,-{Stealth[length=4.5pt,width=4.5pt]}},
  tt collective/.style={draw=ttArrowBlue,line width=0.9pt,-{Stealth[length=5.5pt,width=5.5pt]}},
  tt operator/.style={draw=ttOutlineBlue!76,fill=white,circle,minimum size=0.52cm,
    inner sep=0pt,font=\scriptsize\bfseries,text=black!72},
  tt box operator/.style={draw=ttOutlineBlue!70,fill=white,rounded corners=1.5pt,
    inner xsep=5pt,inner ysep=4pt,font=\scriptsize,text=black!72,align=center},
  tt tensor symbol/.style={font=\scriptsize,text=black!72,anchor=south},
  tt tensor shape/.style={font=\scriptsize,text=ttMutedText,anchor=north,align=center}
}

% Generic numeric matrix. The final argument is a TikZ matrix body using \&
% as the column separator, for example: 1 \& 0 \\ 0 \& 1 \\.
% The matrix itself is a named node, so connectors can use anchors such as
% (q0.east). Labels are separately named as q0-symbol and q0-shape.
\newcommand{\TTMatrix}[7]{%
  \matrix (#1) [matrix of nodes,ampersand replacement=\&,
    nodes in empty cells,
    nodes={draw=ttOutlineBlue!76,fill=#4,minimum width=1.02cm,
      minimum height=0.58cm,inner sep=0pt,font=\scriptsize,
      text=black!78,align=center},
    row sep=-\pgflinewidth,column sep=-\pgflinewidth] at (#2,#3) {#7};
  \node[tt tensor symbol] (#1-symbol) at ($(#1.north)+(0,0.14)$) {#5};
  \node[tt tensor shape] (#1-shape) at ($(#1.south)+(0,-0.15)$) {#6};
}

% Named operator primitive for matmul, add, concat, and other explicit ops.
\newcommand{\TTOperator}[4]{%
  \node[tt operator] (#1) at (#2,#3) {#4};
}

% Named rectangular operator for labels that carry an operand alias or a
% multi-word operation, such as x K_0^T or row-softmax.
\newcommand{\TTBoxOperator}[4]{%
  \node[tt box operator] (#1) at (#2,#3) {#4};
}

% Named waypoint. Raw coordinates belong here, never at connector endpoints.
\newcommand{\TTWaypoint}[3]{%
  \coordinate (#1) at (#2,#3);
}
